\documentclass{amsart}
\usepackage{amsmath,amssymb,amsthm}

\title{Algebraic Derivation of $T^* = 5/7$}
\author{Brayden Ross Sanders / 7Site LLC}
\date{2026}

\newtheorem{theorem}{Theorem}
\newtheorem{corollary}[theorem]{Corollary}

\begin{document}
\maketitle

\begin{abstract}
We prove that the coherence threshold $T^* = 5/7$ is algebraically forced by
the ring $\mathbb{Z}/10\mathbb{Z}$ via four independent derivations:
CRT unit fraction, TSML measurement, generator convergence,
and complement-equivariant fixed point.
\end{abstract}

\section{Four Independent Derivations of $T^* = 5/7$}

\begin{theorem}[$T^*$ Algebraic Identity]
$T^* = 5/7$ is the unique ratio $\mathrm{CREATE}/\mathrm{HARMONY} = 5/7$
in $\mathbb{Z}/10\mathbb{Z}$, established by four independent chains:

\begin{enumerate}
\item[\emph{(a)}] \emph{CRT unit fraction.} At $b=35=5\times7$ and $k=7$:
$|C(7,35)| = |\{1,2,3,4,6\}| = 5$ out of $\{1,\ldots,7\}$, so unit fraction $= 5/7$.

\item[\emph{(b)}] \emph{TSML measurement.} $\mathrm{TSML}[v][j] = \mathrm{HARMONY}=7$
for all non-VOID $v$ and $j \neq \mathrm{VOID}$.
The measurement $M(v) = \mathrm{TSML}[v][v_{\text{journey}}] = 7$ for all non-VOID $v$.
$T^* = \mathrm{destination}/\mathrm{measurement} = 5/7$.

\item[\emph{(c)}] \emph{Generator convergence.}
$\mathrm{CREATE}=5$ is the centroid of $(\mathbb{Z}/10\mathbb{Z})^* = \{1,3,7,9\}$:
$(1+3+7+9)/4 = 5$.
$\mathrm{HARMONY}=7 = g^3 = g^{-1} \pmod{10}$ for $g=3$
(the unique primitive root compatible with $T^* \in (0,1)$).

\item[\emph{(d)}] \emph{Complement-equivariant fixed point.}
$\mathrm{CREATE}=5$ is the unique complement-equivariant ODD-output fixed point:
$2F(5) \equiv 0 \pmod{10}$ forces $F(5) \in \{0,5\}$; since $F(5) \in \mathrm{ODD}$, $F(5)=5$.
\end{enumerate}
\end{theorem}

\begin{theorem}[Generator Uniqueness]
$g=3$ is the unique primitive root of $(\mathbb{Z}/10\mathbb{Z})^*$ such that $T^* \in (0,1)$.
Under $g=7$: $T^* = 5/3 > 1$ (inadmissible).
\end{theorem}

\begin{proof}
Primitive roots of $(\mathbb{Z}/10\mathbb{Z})^*$ are $g \in \{3,7\}$ (order 4 elements).
Under $g=3$: $\mathrm{HARMONY} = g^3 = 27 \bmod 10 = 7$, $T^* = 5/7 < 1$. \checkmark
Under $g=7$: $\mathrm{HARMONY} = g^3 = 343 \bmod 10 = 3$, $T^* = 5/3 > 1$. \texttimes
\end{proof}

\begin{corollary}[Corridor Portrait]
The four spine-forced corridor positions satisfy the strict ordering:
\[
  W < \tfrac{1}{2} < \tfrac{7}{10} < T^* < 1,
  \quad \text{i.e.,} \quad
  \tfrac{3}{50} < \tfrac{1}{2} < \tfrac{7}{10} < \tfrac{5}{7} < 1,
\]
with strict amplitude reversal ($\mathrm{sinc}^2$ is strictly decreasing on $(0,1)$):
\[
  \mathrm{sinc}^2(W) > \mathrm{sinc}^2(\tfrac{1}{2}) >
  \mathrm{sinc}^2(\tfrac{7}{10}) > \mathrm{sinc}^2(T^*).
\]
The inheritance boundary at $t=1/2$ separates ring-forced territory ($t < 1/2$)
from generator-forced territory ($t > 1/2$).
\end{corollary}

\end{document}
